\documentclass[12pt,a4paper]{ctexart}

% === 页面设置 ===
\usepackage[top=2.5cm,bottom=2.5cm,left=2.5cm,right=2.5cm]{geometry}

% === 数学和物理 ===
\usepackage{amsmath,amssymb,amsthm}
\usepackage{physics}

% === 图形和表格 ===
\usepackage{graphicx}
\usepackage{float}
\usepackage{booktabs}
\usepackage{array}
\usepackage{caption}
\usepackage{subcaption}

% === 代码 ===
\usepackage{listings}
\usepackage{xcolor}
\usepackage{upquote}

\definecolor{codebg}{rgb}{0.95,0.95,0.95}
\definecolor{codeframe}{rgb}{0.7,0.7,0.7}
\definecolor{commentcolor}{rgb}{0.3,0.5,0.3}
\definecolor{stringcolor}{rgb}{0.7,0.3,0.1}
\definecolor{keywordcolor}{rgb}{0.2,0.2,0.7}

\lstset{
    basicstyle=\ttfamily\small,
    backgroundcolor=\color{codebg},
    frame=single,
    rulecolor=\color{codeframe},
    breaklines=true,
    numbers=left,
    numberstyle=\tiny\color{gray},
    xleftmargin=2em,
    framexleftmargin=1.5em,
    commentstyle=\color{commentcolor},
    stringstyle=\color{stringcolor},
    keywordstyle=\color{keywordcolor}\bfseries,
    showstringspaces=false,
    tabsize=4,
    captionpos=b,
    language=Python,
    morekeywords={self,True,False,None},
    literate={÷}{{$\div$}}1 {≤}{{$\leq$}}1 {≥}{{$\geq$}}1 {≠}{{$\neq$}}1
           {→}{{$\rightarrow$}}2 {←}{{$\leftarrow$}}2 {λ}{{$\lambda$}}1
}

% === 引用 ===
\usepackage{hyperref}
\usepackage[numbers,sort&compress]{natbib}

% === 其他 ===
\usepackage{enumitem}
\usepackage{fancyhdr}
\usepackage{multirow}

\hypersetup{
    colorlinks=true,
    linkcolor=blue,
    citecolor=blue,
    urlcolor=blue,
}

% === 标题信息 ===
\title{\heiti LQCD Master \\[3pt] \large 格点QCD自主AI科学家框架使用手册}
\author{张鑫}
\date{\today}

\begin{document}

\maketitle

\begin{abstract}
\noindent
LQCD Master是一个面向格点量子色动力学（Lattice QCD）的自主AI科学家框架。它将自然语言描述的物理任务自动翻译为可执行的PyQUDA GPU计算程序及Slurm集群提交脚本，并通过Planner--Executor双阶段流水线实现从物理需求到集群计算的端到端自动化。本手册详细介绍LQCD Master的架构设计、安装配置、命令行使用、技能系统、基准测试以及核心API接口。
\end{abstract}

\tableofcontents
\newpage

%=============================================================================
\section{项目概述}
%=============================================================================

LQCD Master的核心目标是将自然语言描述的格点QCD物理请求自动转换为可执行的PyQUDA程序和Slurm提交脚本。系统结合了\textbf{Planner}（规划器）和\textbf{Executor}（执行器）两个智能体，形成一条可复现、有人类监督的计算流水线。

\subsection{主要特性}
\begin{itemize}[leftmargin=*]
    \item \textbf{自然语言接口}：用户只需用中文或英文描述物理任务（如"计算π介子两点关联函数"），系统自动生成计算方案和代码
    \item \textbf{双阶段流水线}：Planner负责物理方案生成（含评审--重写循环），Executor负责代码生成（含静态分析和自动调试）
    \item \textbf{人在回路}：每个阶段都设有人工检查点，系统不会在未经确认的情况下向集群提交作业
    \item \textbf{自动调试}：若静态分析检测到生成代码中的错误，Executor会自动重写（上限为\texttt{executor\_static\_check\_rounds}轮）
    \item \textbf{技能路由}：领域知识（关联函数理论、PyQUDA API、分析方法）根据当前任务相关性注入LLM提示词
    \item \textbf{全面基准测试}：70个独立验证任务，覆盖5类可观测量
\end{itemize}

\subsection{流水线概览}
\begin{figure}[H]
\centering
\fbox{\parbox{0.9\textwidth}{\centering
\vspace{0.5cm}
\texttt{自然语言任务} $\rightarrow$ \textbf{Planner} (方案生成+评审+重写) $\rightarrow$ [人工检查点] $\rightarrow$ \\
\textbf{Executor} (代码生成+自动调试) $\rightarrow$ [人工检查点] $\rightarrow$ \textbf{Slurm提交} $\rightarrow$ GPU计算结果
\vspace{0.5cm}}}
\caption{LQCD Master 流水线示意图}
\label{fig:pipeline}
\end{figure}

%=============================================================================
\section{项目结构}
%=============================================================================

\begin{lstlisting}[caption={LQCD Master 项目目录结构}, label={lst:structure}]
.
├── run.py                       # CLI入口点
├── requirements.txt             # Python依赖
├── .env.template                # 环境变量模板
├── core_architecture/           # Planner → Executor → Submit 流水线
│   ├── orchestrator.py          # 工作流编排器
│   ├── planner.py               # 规划器Agent
│   └── executor.py              # 执行器Agent
├── utils/                       # LLM客户端、工具、I/O、技能系统
│   ├── llm_client.py            # OpenAI兼容LLM客户端
│   ├── tool_client.py           # 内置工具客户端（搜索/解析）
│   ├── submit_tool.py           # Slurm提交工具
│   ├── config_utils.py          # 配置加载工具
│   ├── io_utils.py              # 文件I/O工具
│   ├── prompt_loader.py         # 提示词加载器
│   ├── skill_utils.py           # 技能系统核心
│   ├── static_check.py          # 静态代码分析
│   └── generate_einsum/         # Wick收缩与einsum生成
├── prompts/                     # 各阶段提示词模板
│   ├── planner/                 # solve / critique / rewrite
│   └── executor/                # generate / critique / rewrite
├── configs/                     # 集群、系综、技能配置
│   ├── config.yaml              # Slurm/Runtime配置
│   ├── ensemble_presets.yaml    # 系综预设
│   └── skills.yaml              # 技能路由配置
├── skills/                      # 领域知识模块
│   ├── lqcd-analysis/           # 数据分析技能
│   ├── lqcd-physics-spectrum/   # 能谱学技能
│   ├── lqcd-physics-correlator/ # 关联函数理论
│   ├── pyquda-tool/             # PyQUDA API指南
│   └── pyquda-gauge/            # 规范场操作
├── benchmark/                   # 70个验证任务 + 标准方法参考数据
└── experiments/                 # 实验结果（按模型分类）
\end{lstlisting}

%=============================================================================
\section{安装}
%=============================================================================

\subsection{环境要求}

\begin{table}[H]
\centering
\caption{运行环境要求}
\label{tab:requirements}
\begin{tabular}{@{}ll@{}}
\toprule
\textbf{组件} & \textbf{要求} \\
\midrule
Python & $\ge$ 3.10 \\
pip依赖 & openai, prompt\_toolkit, PyYAML, numpy, opt\_einsum \\
GPU后端 & CuPy（CUDA或ROCm） \\
格点QCD & PyQUDA + QUDA（从源码编译安装） \\
集群 & Slurm + GPU分区 \\
外部API & OpenAI兼容LLM端点 + Serper.dev（网页搜索） \\
\bottomrule
\end{tabular}
\end{table}

\subsection{安装步骤}

\begin{lstlisting}[caption={LQCD Master 安装流程}, label={lst:install}]
# 1. 克隆仓库
git clone https://github.com/sjtu-sai-agents/LQCD_Master.git
cd LQCD_Master

# 2. 安装Python依赖
pip install -r requirements.txt

# 3. 配置环境变量
cp .env.template .env
# 编辑 .env 文件，填入LLM API密钥和Serper密钥
\end{lstlisting}

\subsection{环境配置（\texttt{.env}）}

\begin{lstlisting}[caption={.env 配置文件示例}, label={lst:env}]
OPENAI_API_KEY=sk-...         # LLM API密钥（必需）
OPENAI_BASE_URL=...           # API端点（默认: api.deepseek.com/v1）
OPENAI_MODEL=...              # 模型名称（默认: deepseek-v4-pro）
SERPER_API_KEY=...            # Serper.dev密钥（网页搜索工具）
\end{lstlisting}

%=============================================================================
\section{配置}
%=============================================================================

\subsection{集群与系综配置（\texttt{configs/}）}

\texttt{configs/config.yaml}定义了Slurm分区、模块加载、MPI启动和QUDA环境路径。\texttt{configs/ensemble\_presets.yaml}定义了7个格点系综的预设信息，包括规范组态路径、夸克质量、clover系数和多重网格参数。

\subsubsection{配置文件结构}

\begin{lstlisting}[caption={configs/config.yaml 结构}, label={lst:config}]
script:     # 脚本模板设置
  shell: /bin/bash
slurm:      # Slurm作业调度参数
  partition: gpu
  nodes: 1
  gres: gpu:4
  ntasks_per_node: 4
  ntasks_per_socket: 4
runtime:    # 运行时环境
  modules: [cuda/12.0, openmpi/4.1]
  exports: {QUDA_PATH: /path/to/quda}
  activate: conda activate lqcd
  source: [source ./env.sh]
  launch: mpirun -np
ensemble:   # 系综配置
  preset: C24P29
  cfg_num: [10000, 10100, 10200]
  process_grid: [2, 2, 1, 1]
measurement:  # 测量参数
solver:       # 求解器参数
output:       # 输出设置
workflow:     # 工作流设置
  executor_static_check_rounds: 5
  executor_dynamic_test_rounds: 2
\end{lstlisting}

\subsection{技能路由配置（\texttt{configs/skills.yaml}）}

控制哪些领域知识模块在每个流水线阶段被注入。系统内置了五个技能模块：
\begin{itemize}[leftmargin=*]
    \item \textbf{lqcd-physics-correlator}：关联函数理论（π介子、质子等强子谱学）
    \item \textbf{lqcd-physics-spectrum}：能谱分析
    \item \textbf{lqcd-analysis}：数据分析方法
    \item \textbf{pyquda-tool}：PyQUDA API使用指南
    \item \textbf{pyquda-gauge}：规范场操作
\end{itemize}

%=============================================================================
\section{核心架构}
%=============================================================================

\subsection{WorkflowOrchestrator（工作流编排器）}

\texttt{WorkflowOrchestrator}是整个系统的主控制器，负责协调Planner和Executor两个子智能体，管理人工检查点，以及处理Slurm作业的提交和监控。

\begin{lstlisting}[caption={WorkflowOrchestrator 使用示例}, label={lst:orchestrator}]
from core_architecture.orchestrator import WorkflowOrchestrator
from utils.llm_client import LQCDLLMClient
from utils.submit_tool import SlurmSubmitTool
from utils.tool_client import BuiltinToolClient

llm = LQCDLLMClient(dotenv_path=".env")
tool_client = BuiltinToolClient(dotenv_path=".env")
submit_tool = SlurmSubmitTool()

orchestrator = WorkflowOrchestrator(
    task="Calculate the pion two-point correlator",
    run_dir=Path("runs/demo"),
    llm_client=llm,
    tool_client=tool_client,
    submit_tool=submit_tool,
    test_mode=True,
    non_interactive=False,
)
final_state = orchestrator.run()
\end{lstlisting}

\subsection{PlannerAgent（规划器）}

Planner通过三阶段流程将自然语言任务转化为结构化的物理计算方案：

\begin{enumerate}[leftmargin=*]
    \item \textbf{Solve}：根据任务描述和系综信息生成初始计算方案（YAML + Markdown）
    \item \textbf{Critique}：对方案进行评审，识别问题和风险，提出修订建议
    \item \textbf{Rewrite}：根据评审意见修订方案
\end{enumerate}

输出格式包含：
\begin{itemize}[leftmargin=*]
    \item \texttt{plan\_yaml}：结构化计算方案（含任务、系综、物理、测量、输出五个部分）
    \item \texttt{summary\_md}：人类可读的方案摘要
    \item \texttt{citations}：验证过的引用URL列表
\end{itemize}

\subsection{ExecutorAgent（执行器）}

Executor将Planner输出的计算方案转换为可执行的Python代码和Slurm脚本：

\begin{enumerate}[leftmargin=*]
    \item \textbf{Generate}：生成PyQUDA主程序 + 测试提交脚本 + 完整提交脚本
    \item \textbf{Static Analysis}：静态分析生成的代码，检测语法错误和导入问题
    \item \textbf{Critique}：结合静态分析结果评审代码质量
    \item \textbf{Rewrite}：根据评审意见和静态分析错误自动重写（最多\texttt{executor\_static\_check\_rounds}轮）
\end{enumerate}

生成产物：
\begin{itemize}[leftmargin=*]
    \item \texttt{main.py}：PyQUDA计算主程序
    \item \texttt{submit\_test.sh}：测试提交脚本（小规模验证运行）
    \item \texttt{submit\_full.sh}：完整提交脚本（生产运行）
    \item \texttt{static\_analysis.json}：静态分析报告
    \item \texttt{critique.json}：代码评审报告
\end{itemize}

\subsection{SkillRunner（技能运行器）}

技能系统通过\texttt{SkillRunner}实现领域知识注入：

\begin{itemize}[leftmargin=*]
    \item \textbf{SkillRegistry}：从\texttt{skills/}目录加载技能模块
    \item \textbf{SkillRouter}：根据任务内容匹配最相关的技能
    \item \textbf{SkillMessageAssembler}：将技能内容组装到LLM提示词中
    \item \textbf{ToolRegistry}：注册和管理内置工具（web\_search, web\_parse）
\end{itemize}

%=============================================================================
\section{命令行使用}
%=============================================================================

\subsection{基本用法}

\begin{lstlisting}[caption={LQCD Master 命令行接口}, label={lst:cli}]
# 交互式运行
python run.py

# 指定任务文本
python run.py --task "计算π介子两点关联函数，使用C24P29系综和点源"

# 从文件读取任务
python run.py --task task.txt

# 测试模式（仅生成代码，不提交作业）
python run.py --task "..." --test

# 非交互模式（自动接受所有检查点）
python run.py --task "..." --non-interactive

# 指定运行目录
python run.py --task "..." --run-dir runs/my_experiment

# 列出可用技能
python run.py --list-skills
\end{lstlisting}

\subsection{命令行参数完整列表}

\begin{table}[H]
\centering
\caption{CLI参数说明}
\label{tab:cli_args}
\begin{tabular}{@{}llp{6cm}@{}}
\toprule
\textbf{参数} & \textbf{类型} & \textbf{说明} \\
\midrule
\texttt{--task} & str & 任务描述文本或包含任务描述的文件路径 \\
\texttt{--test} & flag & 测试模式：仅生成代码，跳过Slurm提交 \\
\texttt{--run-dir} & str & 输出目录（默认：\texttt{runs/<timestamp>}） \\
\texttt{--non-interactive} & flag & 非交互模式：自动接受所有检查点 \\
\texttt{--dotenv-path} & str & 自定义\texttt{.env}文件路径 \\
\texttt{--list-skills} & flag & 列出可用技能并退出 \\
\bottomrule
\end{tabular}
\end{table}

%=============================================================================
\section{示例演示}
%=============================================================================

\subsection{示例1：计算π介子两点关联函数}

\begin{lstlisting}[caption={运行π介子两点关联函数计算}]
# 使用非交互模式快速运行
python run.py --task \
  "Calculate the pion two-point correlator on the C24P29 ensemble using a point source at zero momentum." \
  --test --non-interactive
\end{lstlisting}

Planner生成的方案YAML示例：
\begin{lstlisting}[caption={Planner输出的plan.yaml示例}]
task: "Calculate pion two-point correlator"
ensemble:
  preset: C24P29
  cfg_num: [10000]
  Lx: 24, Ly: 24, Lz: 24, Lt: 72
  beta: 6.20
physics:
  observable: pion_mass
  hadron: pion
  interpolator: gamma5
measurement:
  method: distillation
  nev: 100
  mom_list: [[0, 0, 0, 0]]
output:
  format: npy
  path: ./output/
extras: {}
\end{lstlisting}

\subsection{示例2：批量运行脚本}

\begin{lstlisting}[caption={批量运行实验脚本示例}, language=bash]
#!/bin/bash
# 批量计算多种强子的两点关联函数

TASKS=(
  "Calculate pion two-point correlator on C24P29"
  "Calculate kaon two-point correlator on C24P29"
  "Calculate proton two-point correlator on C24P29"
)

for task in "${TASKS[@]}"; do
  python run.py --task "$task" \
    --test --non-interactive \
    --run-dir "runs/batch_$(date +%Y%m%d_%H%M%S)"
done
\end{lstlisting}

%=============================================================================
\section{技能系统}
%=============================================================================

\subsection{技能模块结构}

每个技能模块位于\texttt{skills/}目录下，包含一个\texttt{SKILL.md}文件：

\begin{lstlisting}[caption={技能模块目录结构}]
skills/
├── lqcd-analysis/SKILL.md            # 数据分析指导
├── lqcd-physics-spectrum/SKILL.md    # 强子能谱学
├── lqcd-physics-correlator/SKILL.md   # 关联函数理论
│   └── reference/                    # 参考数据
│       ├── pion_mass.md
│       ├── proton_mass.md
│       └── rho_mass.md
├── pyquda-tool/SKILL.md              # PyQUDA API使用指南
│   └── reference/
│       └── baryon_3pt_code.md
└── pyquda-gauge/SKILL.md             # 规范场操作
    └── reference/
        └── wilson_loop.md
\end{lstlisting}

\subsection{技能路由机制}

技能路由根据以下规则将相关技能注入到各阶段的提示词中：
\begin{itemize}[leftmargin=*]
    \item \textbf{Planner阶段}：选择与观测物类型（2pt、3pt、Wilson loop等）匹配的技能
    \item \textbf{Executor阶段}：选择与代码生成相关的PyQUDA API技能
    \item 技能注入优先级在\texttt{configs/skills.yaml}中配置
\end{itemize}

%=============================================================================
\section{基准测试}
%=============================================================================

\subsection{测试任务概览}

\texttt{benchmark/}目录定义了70个独立验证任务，涵盖5类可观测量，全部在C24P29系综（cfg 10000，点源，零动量）上评估。

\begin{table}[H]
\centering
\caption{基准测试任务分类}
\label{tab:benchmark}
\begin{tabular}{@{}lccp{5cm}@{}}
\toprule
\textbf{可观测量} & \textbf{任务数} & \textbf{示例} \\
\midrule
2pt local & 20 & $\pi$, K, $\rho$, D, D$_s$, J/$\psi$, p, $\Lambda$, $\Xi$, $\Sigma$, $\Lambda_c$, $\Xi_c$, $\Omega_c$ \\
2pt nonlocal & 10 & $\pi$, K, $\rho$, D, D$_s$, J/$\psi$, D*, D$_s*$ (Wilson线位移 z = 0..10) \\
Wilson loop & 12 & W(R,T), R,T $\in$ \{1,2,3,4\}, 平均于XT, YT, ZT平面 \\
3pt meson & 13 & D$\to$K/$\pi$, B$\to$D/K/$\pi$, K$\to\pi$, B$_c\to$J/$\psi$ \\
3pt baryon & 15 & p$\to$p, $\Lambda\to\Lambda$, $\Lambda_c\to\Lambda$, $\Lambda_b\to\Lambda_c$ \\
\bottomrule
\end{tabular}
\end{table}

\subsection{实验结果}

\begin{table}[H]
\centering
\caption{70任务套件上的交叉验证结果}
\label{tab:results}
\begin{tabular}{@{}lccc@{}}
\toprule
\textbf{实验} & \textbf{骨干模型} & \textbf{准确匹配} & \textbf{成功率} \\
\midrule
LQCD Master GPT-5.4 & GPT-5.4 & 63 & 90.0\% \\
LQCD Master DeepSeek & DeepSeek V4 Pro & 56 & 80.0\% \\
\bottomrule
\end{tabular}
\end{table}

注：准确匹配定义为与标准方法在机器精度（$|\Delta| < 10^{-12}$）或相对误差$< 10^{-3}$下一致。

%=============================================================================
\section{API接口清单}
%=============================================================================

\subsection{LQCDLLMClient}

OpenAI兼容的LLM客户端，支持自动加载\texttt{.env}配置。

\begin{lstlisting}[caption={LQCDLLMClient 接口}]
class LQCDLLMClient:
    """OpenAI兼容客户端，自动加载.env配置"""

    def __init__(
        self,
        *,
        api_key: str | None = None,
        base_url: str | None = None,
        model: str | None = None,
        max_retries: int | None = None,
        dotenv_path: str | Path = ".env",
        extra_body: dict[str, Any] | None = None,
    ): ...

    def chat(
        self,
        messages: list[dict[str, str]],
        *,
        stream: bool = False,
        temperature: float | None = None,
        max_tokens: int | None = None,
        model: str | None = None,
    ) -> str: ...

    def chat_json(
        self,
        messages: list[dict[str, str]],
        *,
        stream: bool = False,
        temperature: float | None = None,
        max_tokens: int | None = None,
        model: str | None = None,
    ) -> dict[str, Any]: ...

    def chat_json_with_tools(
        self,
        messages: list[dict[str, Any]],
        *,
        tools: list[dict[str, Any]],
        tool_handler: Any,
        max_turns: int = 8,
        temperature: float | None = None,
        max_tokens: int | None = None,
        model: str | None = None,
    ) -> tuple[dict[str, Any], list[dict[str, Any]]]: ...
\end{lstlisting}

\subsection{WorkflowOrchestrator}

\begin{lstlisting}[caption={WorkflowOrchestrator 接口}]
class WorkflowOrchestrator:
    """工作流编排器 — Planner → Executor → Submit"""

    def __init__(
        self,
        *,
        task: str,
        run_dir: Path,
        llm_client: Any,
        tool_client: Any,
        submit_tool: Any,
        test_mode: bool = False,
        non_interactive: bool = False,
        max_revision_rounds: int = 3,
    ): ...

    def run(self) -> dict[str, Any]: ...
    # 返回完整状态字典：planner, executor, test_submission, full_submission
\end{lstlisting}

\subsection{PlannerAgent}

\begin{lstlisting}[caption={PlannerAgent 接口}]
class PlannerAgent:
    """Planner Agent — 生成物理计算方案"""

    def __init__(
        self,
        llm_client: Any,
        tool_client: Any,
        run_dir: Path,
    ): ...

    def run(self, task: str) -> dict[str, Any]: ...
    # 返回: {plan_yaml, summary_md, citations}

    def rewrite_with_feedback(
        self,
        task: str,
        current: dict[str, Any],
        feedback: str,
        version: int
    ) -> dict[str, Any]: ...

    def validate_plan_coverage(self, plan_yaml: str) -> dict[str, Any]: ...
    # 检查plan.yaml是否覆盖所有必需字段

    def debug_state(self) -> dict[str, Any]: ...
    # 返回内部状态（solve/critique/rewrite payload, stage traces）
\end{lstlisting}

\subsection{ExecutorAgent}

\begin{lstlisting}[caption={ExecutorAgent 接口}]
class ExecutorAgent:
    """Executor Agent — 生成PyQUDA代码和Slurm脚本"""

    def __init__(
        self,
        llm_client: Any,
        submit_tool: Any,
        run_dir: Path,
        tool_client: Any | None = None,
        max_static_check_rounds: int | None = None,
    ): ...

    def generate_bundle(
        self,
        *,
        task: str,
        planner_output: dict[str, Any],
        version: int,
        feedback: str | None = None,
        previous_bundle: dict[str, Any] | None = None,
    ) -> dict[str, Any]: ...
    # 返回: {main_program_path, test_submit_path, full_submit_path,
    #         static_analysis_report_path, critique_report_path,
    #         static_analysis, critique}

    def submit_test(self, bundle: dict[str, Any], version: int) -> dict[str, Any]: ...

    def submit_full(self, bundle: dict[str, Any], version: int) -> dict[str, Any]: ...

    def collect_evidence_paths(self, bundle: dict[str, Any], version: int) -> list[str]: ...
\end{lstlisting}

\subsection{BuiltinToolClient}

\begin{lstlisting}[caption={BuiltinToolClient 接口}]
class BuiltinToolClient:
    """内置工具客户端 — web_search 和 web_parse"""

    def __init__(self, dotenv_path: str | Path = ".env"): ...

    def web_search(self, query: str, num: int = 5) -> list[dict[str, Any]]: ...
    # 通过Serper API执行网页搜索

    def web_parse(self, url: str) -> dict[str, Any]: ...
    # 解析网页内容，返回文本和元数据
\end{lstlisting}

\subsection{SlurmSubmitTool}

\begin{lstlisting}[caption={SlurmSubmitTool 接口}]
class SlurmSubmitTool:
    """Slurm作业提交工具"""

    def submit(self, script_path: str) -> dict[str, Any]: ...
    # 通过sbatch提交作业脚本
    # 返回: {ok, job_id, stderr}
\end{lstlisting}

\subsection{SkillRegistry / SkillRouter / SkillRunner}

\begin{lstlisting}[caption={技能系统接口}]
class SkillRegistry:
    """技能注册表 — 从skills/目录加载"""
    def __init__(self, skills_dir: Path): ...
    def load(self) -> None: ...
    def names(self) -> list[str]: ...
    def get(self, name: str) -> Skill | None: ...

class SkillRouter:
    """技能路由器 — 根据任务匹配最相关技能"""
    def __init__(self, registry: SkillRegistry, config_path: Path): ...
    def route(self, context: SkillContext) -> list[str]: ...

class SkillRunner:
    """技能运行器 — 执行技能并管理LLM交互"""
    def __init__(
        self,
        llm_client: Any,
        registry: SkillRegistry,
        router: SkillRouter,
        assembler: SkillMessageAssembler,
        tool_registry: ToolRegistry,
    ): ...
    def run(
        self,
        context: SkillContext,
        system_prompt: str,
        user_message: str,
        fallback_tools: list[str] | None = None,
    ) -> tuple[dict[str, Any], list[dict[str, Any]], dict[str, Any]]: ...
\end{lstlisting}

\subsection{辅助工具模块}

\begin{lstlisting}[caption={辅助工具函数}]
# utils/prompt_loader.py
class PromptLoader:
    """提示词加载器"""
    def __init__(self, prompts_dir: Path): ...
    def load(self, relative_path: str) -> str: ...

# utils/config_utils.py
def load_yaml_mapping(path: Path) -> dict[str, Any]: ...
def resolve_submit_config_presets(
    config: dict[str, Any], presets: dict[str, Any]
) -> dict[str, Any]: ...
def build_prompt_ensemble_yaml(
    config: dict[str, Any], presets: dict[str, Any]
) -> str: ...
def build_prompt_submit_config_yaml(
    config: dict[str, Any], presets: dict[str, Any]
) -> str: ...
def build_runtime_launch(runtime_cfg: dict[str, Any]) -> str: ...

# utils/io_utils.py
def read_text(path: Path) -> str: ...
def write_text(path: Path, content: str) -> None: ...
def write_json(path: Path, data: Any) -> None: ...
def extract_json_from_text(text: str) -> dict[str, Any]: ...

# utils/static_check.py
def analyze_bundle_static(
    *,
    main_program: str,
    test_submit_script: str,
    full_submit_script: str,
) -> dict[str, Any]: ...
# 对生成的代码和脚本进行静态分析
# 返回: {main_errors, test_sh_errors, full_sh_errors}
\end{lstlisting}

\subsection{Wick收缩生成器}

\begin{lstlisting}[caption={Wick收缩与einsum生成}]
# utils/generate_einsum/
from utils.generate_einsum import hadron_operator, _wick_translate

# hadron_operator.py — 强子插值算符定义
class HadronOperator: ...

# _wick_translate.py — Wick定理展开与einsum代码生成
def generate_contraction_einsum(
    operator: HadronOperator, ...
) -> str: ...
\end{lstlisting}

%=============================================================================
\section{依赖关系}
%=============================================================================

LQCD Master依赖于以下组件：

\begin{table}[H]
\centering
\caption{依赖关系}
\label{tab:deps}
\begin{tabular}{@{}lll@{}}
\toprule
\textbf{依赖} & \textbf{版本要求} & \textbf{说明} \\
\midrule
Python & $\ge$ 3.10 & 核心语言 \\
openai & $\ge$ 1.0 & OpenAI兼容API客户端 \\
prompt\_toolkit & $\ge$ 3.0 & 交互式命令行提示 \\
PyYAML & $\ge$ 6.0 & YAML配置解析 \\
numpy & $\ge$ 1.24 & 数值计算 \\
opt\_einsum & $\ge$ 3.0 & einsum优化 \\
cupy & $\ge$ 12.0 & GPU加速（CUDA/ROCm） \\
PyQUDA & 最新 & Python-QUDA接口 \\
QUDA & develop分支 & GPU格点QCD库 \\
Slurm & — & 集群作业调度 \\
\bottomrule
\end{tabular}
\end{table}

%=============================================================================
\section{开发与贡献}
%=============================================================================

\subsection{自定义集群配置}

\begin{enumerate}[leftmargin=*]
    \item 编辑\texttt{configs/config.yaml}：设置Slurm分区、模块加载和MPI启动参数
    \item 编辑\texttt{configs/ensemble\_presets.yaml}：添加你的系综信息（规范组态路径、夸克质量等）
\end{enumerate}

\subsection{添加新技能}

\begin{enumerate}[leftmargin=*]
    \item 在\texttt{skills/}下创建新目录
    \item 编写\texttt{SKILL.md}文件（领域知识和LLM指令）
    \item 在\texttt{configs/skills.yaml}中注册技能并配置路由规则
\end{enumerate}

\subsection{扩展基准测试}

\begin{enumerate}[leftmargin=*]
    \item 在\texttt{benchmark/}中添加新任务定义
    \item 提供手写的标准方法PyQUDA参考实现及数值输出
    \item 运行交叉验证以评估新LLM骨干模型
\end{enumerate}

\end{document}
